\section{Where to perturb a ViT}
\label{sec:results}

This section walks from the original placement to PSBD-TM in 4 moves, each on the same backdoored ViT-B/16 models. We first apply dropout to the residual stream, as our reading of the original placement. We then keep dropout and move it to the other sites of the block. We then keep the best site and change what is injected there. Last, we restrict the winning placement to a band of blocks. Every number below is read at the rate the original rule assigns, and the per-model numbers behind each table are in \cref{app:results}.

\subsection{Dropout on the residual stream}
\label{sec:results:residual}

A ViT block has 2 residual adds, so PSBD-RD applies dropout to the stream after both of them. \Cref{tab:staircase-residual} reads it beside its nearest variants, dropout after the attention add only, dropout on the branch outputs just before the adds, and the latter restricted to a third of the network.

% generated by scripts/paper/tab_staircase.py from results/coverage/coverage.json, results/<folder>/psbd_metrics.json (65 models) at 2026-09-11T04:17:09.248134+00:00, commit 4a04e5dd204b14d8368fb7eeef7092bd331e00d9-dirty. Do not edit by hand.
\begin{table*}[htbp]
\centering
\small
\caption{Dropout on the residual stream of ViT-B/16, after both residual adds, after the attention add only, before both adds, and restricted to bands of blocks.}
\label{tab:staircase-residual}
\begin{tabular}{lrrrrl}
\toprule
 & & & \multicolumn{2}{c}{TPR@FPR} & \\
\cmidrule(lr){4-5}
Placement & n & AUROC & 10\% & 20\% & Gain [95\% CI] \\
\midrule
dropout, after both residual adds & 65 & 0.832 & 0.658 & 0.721 & -- \\
dropout, after the attention add only & 9 & 0.816 & 0.534 & 0.665 & -0.017 [-0.039, +0.011] \\
dropout, before both residual adds & 65 & 0.833 & 0.632 & 0.715 & +0.001 [-0.027, +0.027] \\
dropout, before both residual adds, blocks 1 to 4 & 65 & 0.798 & 0.547 & 0.678 & -0.034 [-0.077, +0.009] \\
dropout, before both residual adds, blocks 5 to 8 & 65 & 0.875 & 0.690 & 0.749 & +0.043 [+0.017, +0.072] \\
dropout, before both residual adds, blocks 9 to 12 & 46 & 0.880 & 0.721 & 0.784 & +0.027 [-0.012, +0.069] \\
\bottomrule
\end{tabular}
\end{table*}


PSBD-RD reaches a mean AUROC of \StaircaseResidualPostResidualAuroc{} and flags \StaircaseResidualPostResidualTprAtOneZeroPercent{} of the triggered inputs at a 10\% false-positive budget. The mean hides a split: it reads below chance on \StaircaseResidualPostResidualBelowChance{} of the \StaircaseResidualPostResidualN{} models. On BadNets, Blend and LF at the higher poison rates it is close to perfect, and on TaCT and on most models at \PanelRateLow{} poisoning it is close to a coin or worse, which \cref{tab:results-rate-1} shows model by model. Moving the dropout to the branch outputs before the adds changes nothing on average: \StaircaseResidualPreResidualGain{} with an interval of [\StaircaseResidualPreResidualGainLow{}, \StaircaseResidualPreResidualGainHigh{}] that contains 0. Applying it after the attention add only is slightly worse. The site is also fragile across training seeds: \cref{sec:robustness:seeds} shows that PSBD-RD's AUROC moves by \SeedSdPublished{} between seeds of the same model against \SeedSdRecommended{} for PSBD-TM.

Restricting residual dropout to a band of blocks helps a little. The middle third gains \StaircaseResidualPreResidualBlocksFiveEightGain{} over all blocks and the last third \StaircaseResidualPreResidualBlocksNineOneTwoGain{}, while the first third loses \StaircaseResidualPreResidualBlocksOneFourGain{}. The stream in early blocks carries mostly clean content, so perturbing it there spends the budget on the wrong signal.

\subsection{The same dropout at the attention input}
\label{sec:results:sites}

Dropout itself is not the problem. \Cref{tab:staircase-dropout-sites} keeps the perturbation the original used and moves it to the other sites of \cref{fig:vit-block}, in every block.

% generated by scripts/paper/tab_staircase.py from results/coverage/coverage.json, results/<folder>/psbd_metrics.json (65 models) at 2026-09-11T04:17:15.697257+00:00, commit 4a04e5dd204b14d8368fb7eeef7092bd331e00d9-dirty. Do not edit by hand.
\begin{table*}[htbp]
\centering
\small
\caption{Dropout at each site of the ViT-B/16 block, applied in all 12 blocks.}
\label{tab:staircase-dropout-sites}
\begin{tabular}{lrrrrl}
\toprule
 & & & \multicolumn{2}{c}{TPR@FPR} & \\
\cmidrule(lr){4-5}
Placement & n & AUROC & 10\% & 20\% & Gain [95\% CI] \\
\midrule
dropout, after both residual adds & 65 & 0.832 & 0.658 & 0.721 & -- \\
dropout, before both residual adds & 65 & 0.833 & 0.632 & 0.715 & +0.001 [-0.027, +0.027] \\
dropout, attention input & 37 & 0.925 & 0.792 & 0.897 & +0.039 [-0.000, +0.086] \\
dropout, attention input after norm & 9 & 0.897 & 0.707 & 0.825 & +0.064 [-0.018, +0.157] \\
dropout, MLP input & 9 & 0.812 & 0.507 & 0.641 & -0.022 [-0.053, +0.019] \\
dropout, embedding output & 9 & 0.873 & 0.645 & 0.732 & +0.039 [-0.027, +0.119] \\
\bottomrule
\end{tabular}
\end{table*}


Dropout at the attention input, before the block's LayerNorm, reaches \StaircaseDropoutSitesBeforeAttentionNormAuroc{} against \StaircaseDropoutSitesPostResidualAuroc{} on the residual stream. That is a paired gain of \StaircaseDropoutSitesBeforeAttentionNormGain{} [\StaircaseDropoutSitesBeforeAttentionNormGainLow{}, \StaircaseDropoutSitesBeforeAttentionNormGainHigh{}] and it flags \StaircaseDropoutSitesBeforeAttentionNormTprAtOneZeroPercent{} of the triggered inputs at a 10\% budget against \StaircaseDropoutSitesPostResidualTprAtOneZeroPercent{}. The interval touches 0, so on its own this move is suggestive rather than settled. Dropout at the MLP input is worse than the residual stream at \StaircaseDropoutSitesBeforeMlpNormAuroc{} and dropout on the embedding output sits between the 2 at \StaircaseDropoutSitesAfterEmbeddingAuroc{}. Site matters, but moving dropout alone does not carry the result. What follows does.

\subsection{Masking tokens instead of dropping coordinates}
\label{sec:results:operator}

\Cref{tab:staircase-operators} holds the site at the attention input and changes what is injected there. It then reads the same perturbations at the other sites of the basis, and at the MLP input it reads Gaussian noise on both sides of that sublayer's LayerNorm, which is what separates the operator from the site.

% generated by scripts/paper/tab_staircase.py from results/coverage/coverage.json, results/<folder>/psbd_metrics.json (65 models) at 2026-09-11T04:17:24.565336+00:00, commit 9df2ce36e6044bb46e4438e652a6afdb83e4895c-dirty. Do not edit by hand.
\begin{table*}[htbp]
\centering
\small
\caption{Perturbation operators at the attention input of ViT-B/16, and token masking at the other sites, all 12 blocks.}
\label{tab:staircase-operators}
\begin{tabular}{lrrrrl}
\toprule
 & & & \multicolumn{2}{c}{TPR@FPR} & \\
\cmidrule(lr){4-5}
Placement & n & AUROC & 10\% & 20\% & Gain [95\% CI] \\
\midrule
token mask, attention input & 65 & 0.935 & 0.800 & 0.859 & -- \\
channel mask, attention input & 65 & 0.880 & 0.680 & 0.762 & -0.055 [-0.095, -0.020] \\
gaussian, attention input & 65 & 0.796 & 0.532 & 0.651 & -0.139 [-0.200, -0.083] \\
dropout, attention input & 37 & 0.925 & 0.792 & 0.897 & -0.043 [-0.078, -0.001] \\
token mask, attention output before the add & 65 & 0.936 & 0.812 & 0.890 & +0.001 [-0.025, +0.030] \\
token mask, both sublayer inputs & 65 & 0.903 & 0.724 & 0.798 & -0.033 [-0.062, -0.004] \\
token mask, MLP input & 65 & 0.821 & 0.504 & 0.684 & -0.114 [-0.156, -0.073] \\
gaussian, MLP input after norm & 65 & 0.868 & 0.705 & 0.778 & -0.068 [-0.114, -0.025] \\
token mask, after the attention add & 65 & 0.770 & 0.466 & 0.592 & -0.166 [-0.224, -0.110] \\
channel mask, MLP neurons & 65 & 0.802 & 0.540 & 0.657 & -0.134 [-0.180, -0.087] \\
gain scale, MLP norm output & 65 & 0.896 & 0.773 & 0.833 & -0.039 [-0.072, -0.009] \\
scale up, input pixels & 65 & 0.746 & 0.453 & 0.591 & -0.190 [-0.249, -0.133] \\
dropout, after both residual adds & 65 & 0.832 & 0.658 & 0.721 & -0.103 [-0.160, -0.048] \\
\bottomrule
\end{tabular}
\end{table*}


Masking whole tokens is the best perturbation at that site. It reaches \StaircaseOperatorsBeforeAttentionNormTokenMaskAuroc{} and flags \StaircaseOperatorsBeforeAttentionNormTokenMaskTprAtOneZeroPercent{} of the triggered inputs at a 10\% budget, against \StaircaseOperatorsBeforeAttentionNormChannelMaskAuroc{} for masking whole channels, \StaircaseOperatorsBeforeAttentionNormAuroc{} for dropout and \StaircaseOperatorsBeforeAttentionNormGaussianAuroc{} for Gaussian noise. This is PSBD-TM, and it reads below chance on \StaircaseOperatorsBeforeAttentionNormTokenMaskBelowChance{} of \StaircaseOperatorsBeforeAttentionNormTokenMaskN{} models against \StaircaseOperatorsPostResidualBelowChance{} for PSBD-RD.

Token masking on the attention branch output just before its residual add is its twin. It reads \StaircaseOperatorsBeforeAttentionResidualTokenMaskAuroc{}, a paired difference of \StaircaseOperatorsBeforeAttentionResidualTokenMaskGain{} [\StaircaseOperatorsBeforeAttentionResidualTokenMaskGainLow{}, \StaircaseOperatorsBeforeAttentionResidualTokenMaskGainHigh{}] that we cannot separate from 0, so the 2 sites either side of the attention block behave alike on ViT. We name the attention input for 3 reasons. It is the site whose mechanism \cref{sec:mechanism} can state. It is the more stable of the 2 across training seeds. And on Swin-S the 2 are not close, where the attention input leads its twin by \SwinGainRecommendedMinusTwin{} [\SwinGainRecommendedMinusTwinLow{}, \SwinGainRecommendedMinusTwinHigh{}] (\cref{sec:swin}). Token masking on the stream after the attention add, by contrast, falls to \StaircaseOperatorsAfterAttentionResidualTokenMaskAuroc{}. The 3 sites differ only in whether the residual stream itself is masked, and \cref{sec:mechanism} explains why that matters.

What the perturbation is worth depends on which side of a LayerNorm it lands on rather than on which site it lands at. Injected before a LayerNorm, Gaussian noise loses to token masking at both sites, by \StaircaseOperatorsBeforeAttentionNormGaussianGain{} at the attention input and by \StaircaseOperatorsBeforeMlpNormGaussianGain{} at the MLP input. Injected after the MLP's LayerNorm the same noise reads \StaircaseOperatorsBeforeMlpGaussianAuroc{}, above token masking injected before it, which is the comparison an earlier draft of this paper read as a reversal between 2 adjacent sites. It is not a site effect. Holding the side of the normalization fixed leaves an interaction between site and operator that we cannot separate from 0, and the operator ranking is the same at both sites. \Cref{sec:mechanism:layernorm} traces it to the LayerNorm behind the attention input, which absorbs additive noise but not masking.

Why should whole tokens beat coordinates? Dropout zeroes individual coordinates of the perturbed tensor, and a coordinate is 1 of the \VitWidth{} numbers describing a token. It is renormalized by the next LayerNorm and averaged by the next attention, so dropping a share of coordinates leaves every token present with a slightly noisier description. A token is the unit that attention routes, so masking it removes that token from every attention pattern of the block. \Cref{sec:mechanism} shows that a backdoor is carried by the trigger's own tokens for most of the network and that its direction is not aligned with any coordinate axis, so a coordinate-level perturbation cannot remove the shortcut while a token-level one can. Channel masking, which removes whole coordinates, trails token masking by \StaircaseOperatorsBeforeAttentionNormChannelMaskGain{} at the same site, and on Swin dropout at the same site reads \SwinSameSiteDropoutAuroc{} against \SwinRecommendedAurocAdaptive{} for token masking (\cref{sec:swin}).

\subsection{Is it the family, or the operator?}
\label{sec:results:family}

The basis groups sites into an input-side family and a residual-adjacent family, and the ranking makes the input-side family look like the winner. It is not a property of the families. Pairing within model and holding the operator fixed, the gap between them depends on which operator is held.

With token masking held fixed the gap is \FamilyGapAdaptive{} [\FamilyGapAdaptiveLow{}, \FamilyGapAdaptiveHigh{}] at the deployable rule and \FamilyGapMatched{} [\FamilyGapMatchedLow{}, \FamilyGapMatchedHigh{}] at matched disturbance, where it cannot be separated from 0. With dropout held fixed instead the gap reverses. At matched disturbance it is \FamilyGapDropoutMatched{} [\FamilyGapDropoutMatchedLow{}, \FamilyGapDropoutMatchedHigh{}], which excludes 0 on the other side. At the deployable rule it is \FamilyGapDropoutAdaptive{} [\FamilyGapDropoutAdaptiveLow{}, \FamilyGapDropoutAdaptiveHigh{}]. With dropout the residual-adjacent family is the better one.

A split whose sign depends on the operator is not a fact about sites. Nor do the families group placements that behave alike: per-sample PSU agrees at \AgreementAcrossFamiliesClean{} across the 2 families and at \AgreementWithinInputSideClean{} within the input-side family (\cref{tab:mech-operator-agreement-family}). It is the operator effect of \cref{sec:mechanism:layernorm} read through a grouping that does not hold the operator fixed, which is also what the basis file's own family tags do: under those tags, which mix operators inside each family, the gap is \FamilyGapConfoundedMatched{} [\FamilyGapConfoundedMatchedLow{}, \FamilyGapConfoundedMatchedHigh{}].

% generated by scripts/paper/tab_family_split.py from results/coverage/coverage.json, configs/psbd_basis.json, results/<folder>/psbd_metrics.json (71 cells) at 2026-09-23T18:25:02.156153+00:00, commit 5a1f1c5814dc5676c093f3021490a849712a13f4-dirty. Do not edit by hand.
\begin{table}[htbp]
\centering
\small
\caption{The within-model family gap with the operator held fixed at token masking. Input-side (attention input, MLP input and both sublayer inputs) minus residual-adjacent (attention output before the add and stream after the attention add), each family's per-model value the mean AUROC at the headline quantile over whichever of its members are cached. Read at the matched 0.6 rate and separately at the adaptive 0.8 rate, 5000-resample bootstrap 95\% intervals.}
\label{tab:family-split}
\begin{adjustbox}{max width=\linewidth}
\begin{tabular}{lrrlrrl}
\toprule
subset & n match & gap matched06 & 95\% CI & n adapt & gap adaptive08 & 95\% CI \\
\midrule
all & 71 & +0.016 & [$-$0.001, +0.036] & 71 & +0.045 & [+0.016, +0.075] \\
CIFAR-10 & 18 & +0.027 & [$-$0.008, +0.062] & 18 & +0.005 & [$-$0.055, +0.068] \\
CIFAR-100 & 15 & $-$0.002 & [$-$0.030, +0.028] & 15 & +0.065 & [+0.018, +0.117] \\
GTSRB & 15 & +0.003 & [$-$0.029, +0.050] & 15 & +0.027 & [$-$0.009, +0.068] \\
Tiny ImageNet & 17 & $-$0.008 & [$-$0.026, +0.013] & 17 & +0.041 & [$-$0.013, +0.107] \\
SVHN & 3 & +0.126 & [$-$0.025, +0.239] & 3 & +0.122 & [$-$0.055, +0.267] \\
EuroSAT & 3 & +0.145 & [+0.100, +0.217] & 3 & +0.218 & [+0.033, +0.384] \\
rate 0.01 & 19 & +0.007 & [$-$0.021, +0.037] & 19 & +0.021 & [$-$0.015, +0.062] \\
rate 0.05 & 24 & +0.007 & [$-$0.018, +0.033] & 24 & +0.059 & [+0.025, +0.093] \\
rate 0.1 & 28 & +0.031 & [$-$0.002, +0.067] & 28 & +0.049 & [$-$0.011, +0.111] \\
hard attacks & 33 & +0.043 & [+0.013, +0.075] & 33 & +0.091 & [+0.038, +0.142] \\
easy attacks & 38 & $-$0.007 & [$-$0.025, +0.011] & 38 & +0.005 & [$-$0.017, +0.030] \\
\bottomrule
\end{tabular}
\end{adjustbox}
\end{table}

% generated by scripts/paper/tab_family_split.py from results/coverage/coverage.json, configs/psbd_basis.json, results/<folder>/psbd_metrics.json (71 cells) at 2026-09-23T18:25:11.250396+00:00, commit 5a1f1c5814dc5676c093f3021490a849712a13f4-dirty. Do not edit by hand.
\begin{table}[htbp]
\centering
\small
\caption{The within-model family gap with the operator held fixed at dropout. Input-side (attention input, MLP input and embedding output) minus residual-adjacent (attention output before the add, stream after the attention add, before both residual adds and after both residual adds), each family's per-model value the mean AUROC at the headline quantile over whichever of its members are cached. Read at the matched 0.6 rate and separately at the adaptive 0.8 rate, 5000-resample bootstrap 95\% intervals.}
\label{tab:family-split-dropout}
\begin{adjustbox}{max width=\linewidth}
\begin{tabular}{lrrlrrl}
\toprule
subset & n match & gap matched06 & 95\% CI & n adapt & gap adaptive08 & 95\% CI \\
\midrule
all & 65 & $-$0.066 & [$-$0.085, $-$0.046] & 65 & $-$0.011 & [$-$0.030, +0.009] \\
CIFAR-10 & 18 & $-$0.023 & [$-$0.055, +0.007] & 18 & +0.009 & [$-$0.027, +0.042] \\
CIFAR-100 & 15 & $-$0.067 & [$-$0.104, $-$0.024] & 15 & +0.009 & [$-$0.026, +0.048] \\
GTSRB & 15 & $-$0.079 & [$-$0.119, $-$0.046] & 15 & $-$0.006 & [$-$0.050, +0.051] \\
Tiny ImageNet & 17 & $-$0.096 & [$-$0.135, $-$0.056] & 17 & $-$0.054 & [$-$0.078, $-$0.034] \\
SVHN & 0 & -- & -- & 0 & -- & -- \\
EuroSAT & 0 & -- & -- & 0 & -- & -- \\
rate 0.01 & 19 & $-$0.048 & [$-$0.084, $-$0.008] & 19 & +0.004 & [$-$0.031, +0.043] \\
rate 0.05 & 22 & $-$0.066 & [$-$0.094, $-$0.037] & 22 & $-$0.009 & [$-$0.033, +0.018] \\
rate 0.1 & 24 & $-$0.080 & [$-$0.113, $-$0.045] & 24 & $-$0.024 & [$-$0.058, +0.013] \\
hard attacks & 29 & $-$0.054 & [$-$0.091, $-$0.017] & 29 & $-$0.016 & [$-$0.048, +0.017] \\
easy attacks & 36 & $-$0.075 & [$-$0.093, $-$0.056] & 36 & $-$0.007 & [$-$0.028, +0.017] \\
\bottomrule
\end{tabular}
\end{adjustbox}
\end{table}


\subsection{Every block, not a band}
\label{sec:results:depth}

PSBD-TM perturbs all 12 blocks. \Cref{tab:staircase-bands} restricts it to a third of the network.

% generated by scripts/paper/tab_staircase.py from results/coverage/coverage.json, results/<folder>/psbd_metrics.json (65 models) at 2026-09-11T04:17:30.481601+00:00, commit 9df2ce36e6044bb46e4438e652a6afdb83e4895c-dirty. Do not edit by hand.
\begin{table*}[htbp]
\centering
\small
\caption{Token masking at the attention input of ViT-B/16, restricted to bands of blocks.}
\label{tab:staircase-bands}
\begin{tabular}{lrrrrl}
\toprule
 & & & \multicolumn{2}{c}{TPR@FPR} & \\
\cmidrule(lr){4-5}
Placement & n & AUROC & 10\% & 20\% & Gain [95\% CI] \\
\midrule
token mask, attention input, all 12 blocks & 65 & 0.935 & 0.800 & 0.859 & -- \\
token mask, attention input, blocks 1 to 4 & \pending & \pending & \pending & \pending & \pending \\
token mask, attention input, blocks 5 to 8 & 44 & 0.883 & 0.758 & 0.838 & -0.063 [-0.115, -0.021] \\
token mask, attention input, blocks 9 to 12 & 31 & 0.893 & 0.761 & 0.827 & -0.050 [-0.093, -0.012] \\
\bottomrule
\end{tabular}
\end{table*}


Restricting the mask to blocks 5 to 8 costs \StaircaseBandsBeforeAttentionNormBlocksFiveEightTokenMaskGain{} and to blocks 9 to 12 costs \StaircaseBandsBeforeAttentionNormBlocksNineOneTwoTokenMaskGain{}, so the token mask works best when it acts everywhere. At matched disturbance, where every band is read on the same models (\cref{tab:depth-bands}), the early band is the worst of the 3 at \BandOneFourMinusAllInputSide{} [\BandOneFourMinusAllInputSideLow{}, \BandOneFourMinusAllInputSideHigh{}], which is what the mechanism predicts for a backdoor written in the last third of the network. This is the opposite of what residual dropout showed in \cref{tab:staircase-residual}, where the middle band was the best. \Cref{sec:mechanism} gives the reason. A token mask before attention removes the trigger's tokens at whatever depth they are routed to the class token, so every added block is another chance to remove them, while a perturbation on the stream hits the decision itself and is best kept away from the early blocks.

\subsection{PSBD-TM across attacks, poison rates and datasets}
\label{sec:results:headline}

\Cref{tab:results-rate-1,tab:results-rate-5,tab:results-rate-10} in \cref{app:results} give PSBD-TM and PSBD-RD on every model at \PanelRates{} poisoning, with ASR, clean accuracy, AUROC and TPR@FPR at both budgets. \Cref{tab:gains} summarizes the paired gain by poison rate and dataset.

% generated by scripts/paper/tab_gains.py from results/coverage/coverage.json (71 clearing cells), results/*/psbd_metrics.json at 2026-09-23T19:49:23.033314+00:00, commit 4087ff7b8dddaf7e969f38f703166a0e490aaacf-dirty. Do not edit by hand.
\begin{table}[htbp]
\centering
\small
\caption{Paired AUROC gain of token masking at the attention input over dropout after both residual adds, at the adaptive rule and the headline quantile, with 95\% bootstrap intervals over models from 5000 resamples. Every row is a paired difference on the models carrying both placements.}
\label{tab:gains}
\begin{adjustbox}{max width=\linewidth}
\begin{tabular}{lrrl}
\toprule
subset & models & gain & 95\% interval \\
\midrule
all models & 69 & +0.104 & [+0.052, +0.158] \\
\multicolumn{4}{l}{\emph{by poison rate}} \\
1\% poisoning & 19 & +0.161 & [+0.059, +0.283] \\
5\% poisoning & 23 & +0.093 & [+0.038, +0.158] \\
10\% poisoning & 27 & +0.073 & [$-$0.022, +0.172] \\
\multicolumn{4}{l}{\emph{by dataset}} \\
CIFAR-10 & 18 & +0.083 & [$-$0.058, +0.235] \\
Tiny ImageNet & 17 & +0.049 & [+0.007, +0.105] \\
CIFAR-100 & 15 & +0.150 & [+0.062, +0.245] \\
GTSRB & 15 & +0.142 & [+0.036, +0.277] \\
EuroSAT & 2 & +0.005 & -- \\
SVHN & 2 & +0.234 & -- \\
\bottomrule
\end{tabular}
\end{adjustbox}
\end{table}


The gain grows as the poison rate falls, from \GainsHighestRate{} at \PanelRateHigh{} to \GainsLowestRate{} at \PanelRateLow{}. This is the regime where a defender most needs help, since an attacker who can poison only \PanelRateLow{} of the data plants a weaker shortcut. The gain is positive on \GainsPositiveDatasets{} of \GainsDatasets{} datasets, large on CIFAR-100 and GTSRB and small on Tiny ImageNet, where PSBD-RD is already strong on BadNets, Blend and LF. Removing any single attack leaves a gain of at least \GainsLeaveOneAttackOut{}, with a lower bootstrap bound of \GainsLeaveOneAttackOutLow{}. No 1 attack carries the result. The attack whose removal costs the most is \GainsLeaveOneAttackOutAttack{}. Per attack, PSBD-TM stays at or above \GainsStrongAuroc{} on every dataset and rate only for \GainsStrongAttacks{}. Every other attack has at least 1 model below that, and the floor is \HeadlineFloorAuroc{} on the 2 models \cref{sec:limitations} discusses.

\Cref{fig:psu} shows what PSU looks like on 6 of these models. On the 3 in the top row, a BadNets, a BPP and a Blend model, the triggered inputs keep their prediction under the mask and pile up at a PSU near 0, while the clean inputs lose theirs and pile up near 1, and the threshold at the \HeadlineQuantilePercent{} percentile of the clean validation set sits between the 2 groups with almost nothing on the wrong side. The bottom row shows the 3 weakest readings. On WaNet and SIG on CIFAR-10 the triggered inputs move as much as the clean ones, so the 2 histograms overlap and the threshold cuts both. On TaCT on Tiny ImageNet both groups sit near 1 and only a thin tail of triggered inputs falls below the threshold. \Cref{fig:basis} closes the section with every fully swept placement of the basis on 1 axis.

\begin{figure*}[t]
\centering
\includegraphics[width=\textwidth]{fig_psu_histograms.pdf}
\caption{Histograms of PSU under PSBD-TM for clean and triggered inputs, with the threshold at the 25th percentile of the clean validation set (dashed). Top, 3 models where detection succeeds. Bottom, the 3 weakest readings.}
\label{fig:psu}
\end{figure*}

\begin{figure*}[t]
\centering
\includegraphics[width=0.8\textwidth]{basis_ranking.pdf}
\caption{Mean AUROC of every placement in the basis with a complete sweep, read at the rate the original rule assigns. Placements that mask whole tokens before attention lead. Noise before a LayerNorm and the pixel-space perturbation trail.}
\label{fig:basis}
\end{figure*}
